\section{Introduction}
\label{sec:introduction}
\PARstart{T}{hanks} to its low cost, speed, and accessibility, bone X-ray imaging is widely used in screening, diagnostic support, and treatment monitoring. However, segmenting lesions in this type of image is not simple when the lesion is very small, low-contrast, or partially obscured by the surrounding bone structure. Bounding-box prompts provide an explicit location cue for segmentation. This study examines segmentation conditional on such a cue, using boxes simulated from annotated lesions. The task therefore assumes prior localization of the target region and evaluates the subsequent prediction of its mask.

This is why small-lesion segmentation is rarely just a boundary-drawing problem. When a lesion occupies only a tiny fraction of the image, searching the whole image is both costly and error-prone. The natural approach is to narrow the search space first, then segment within the narrowed region. Fully automatic segmentation methods, typified by U-Net and Attention U-Net \cite{b3,b4}, try to learn both localization and segmentation implicitly within the same network, with no external help at all.

Another approach is to draw on external guidance about the lesion's location. A coarse bounding box around a suspicious region is one such cue, and is not merely an interactive convenience but a practical way to inject spatial knowledge into the model before detailed segmentation begins. Many prompt-based interactive segmentation methods have grown out of this idea, from simple operations such as clicks and scribbles \cite{b5,b6} to recent foundation models such as SAM and SAM-Med2D \cite{b1,b2}.

We investigate a compact architecture in which a box prompt is supplied explicitly to multiple encoder and decoder stages. PGA-UNet represents the box as a Gaussian-smoothed plateau, modulates encoder features through PSG, and adds separately encoded prompt features to decoder skip attention through CAD. This design provides repeated, explicit access to the prompt within the network. Its motivation is to support segmentation when the target occupies a small or visually ambiguous region; preservation of lesion-specific internal features is not measured directly in this study.

The empirical evaluation compares the complete design with conventional ways of incorporating the same box, including a binary input channel and a prompt crop. Ablations examine the observed effects of the prompt representation and conditioning components. Because the conventional channel baseline and the full model use different prompt representations, their comparison evaluates the complete configurations. A binary-prompt PGA-UNet variant provides an additional comparison at a shared prompt representation.

The contributions of this work are as follows.
\begin{itemize}
\item We introduce PGA-UNet, a compact convolutional architecture that incorporates a Gaussian-smoothed box representation into encoder features through a Prompt Spatial Gate and into decoder skip attention through a Conditional Attention Decoder.
\item We evaluate the architecture on BTXRD and FracAtlas using simulated lesion-covering boxes, conventional prompt baselines, and a comparison with SAM-Med2D at matched input resolution and scoring frame.
\item We characterize the effects of the conditioning components, the observed response to a predefined prompt displacement, variability across four additional image-level splits for PGA-UNet, and computational cost in the tested environment.
\end{itemize}

Beyond these contributions, the system also exposes two support signals, computed directly from the trained model's output without ground-truth masks at inference: a self-assessment ranking heuristic, and a short candidate-region shortlist. Both are presented as preliminary and exploratory, and are not central contributions of this article.
